\section{The proposed solution}
\label{sec:algo}

We develop a reverse-sampling construction for the coupon walks of
Section~\ref{sec:model}. A revisit triggers a fresh action, so a single
fixed live edge at each original node does not represent the walk.
We instead construct a finite, randomly layered graph whose backward
reachability probabilities equal the consumption probabilities
$M_{u,s}$. Independent reverse samples then estimate the marginal
spread of every candidate seed, and greedy selection operates directly
on $k$ distinct users. We give the construction, its correctness proof,
and an explicit sampling guarantee below. The time bound accounts for
both the random number of layers and repeated marginal estimation.

\subsection{A Finite Layered Representation of a Coupon Walk}
\label{sec:layered-world}

Write $c_u=p_u^a+p_u^d$ and
$\beta=\min_{u\in V}c_u>0$, using the positive termination probabilities
in Section~\ref{sec:model}. The quantity $\beta$ may depend on the
instance. We introduce an auxiliary coupling of the possible outcomes
from different candidate seeds; its purpose is to enable backward
exploration. It preserves the law of a coupon from each fixed seed.
Different coupons always use independent copies of this construction.

Suppose first that $0<\beta<1$. Draw a nonnegative integer $H$ with
\begin{equation}
\label{eq:random-horizon}
  \Pr[H=h]=\beta(1-\beta)^h,\qquad h=0,1,\ldots.
\end{equation}
At layers $\ell<H$, a state $(u,\ell)$ independently chooses one
residual action with probabilities
\begin{equation}
\label{eq:residual-actions}
\begin{split}
  \widetilde p_u^a&=\frac{p_u^a-\beta p_u^a/c_u}{1-\beta},\qquad
  \widetilde p_u^d=\frac{p_u^d-\beta p_u^d/c_u}{1-\beta},\\
  \widetilde p(u,v)&=\frac{p(u,v)}{1-\beta}.
\end{split}
\end{equation}
These are nonnegative and sum to one because $c_u\geq\beta$ and
Eq.~\eqref{eq:normalize} holds. A transfer leads from $(u,\ell)$
to $(v,\ell+1)$; consumption and dropping terminate the walk.
At layer $H$, every state terminates: node $u$ consumes with probability
$p_u^a/c_u$ and drops with probability $p_u^d/c_u$.
All state actions are independent conditional on $H$.
If $\beta=1$, take $H=0$ and use only this terminal layer; the residual
probabilities in Eq.~\eqref{eq:residual-actions} are then unnecessary.

\begin{lemma}[Preservation of the coupon law]
\label{lem:layered-law}
Starting at $(s,0)$ in the layered construction produces the same
consumption and dropping trajectory distribution as a coupon seeded
at $s$ in Section~\ref{sec:model}.
\end{lemma}

\begin{proof}
Equivalently, draw independent Bernoulli variables of parameter $\beta$
for successive layers and let $H$ be the first successful layer.
At such a layer the coupon is forced to terminate, with the
consumption/drop split above; otherwise it uses a residual action.
Before the action at a visited state, the current Bernoulli variable
is independent of the earlier path. For $\beta<1$, the resulting
probabilities of consumption, dropping, and transfer to $v$ are
\begin{align*}
  \beta\frac{p_u^a}{c_u}+(1-\beta)\widetilde p_u^a&=p_u^a,\\
  \beta\frac{p_u^d}{c_u}+(1-\beta)\widetilde p_u^d&=p_u^d,\\
  (1-\beta)\widetilde p(u,v)&=p(u,v).
\end{align*}
Every subsequent step uses a new layer and fresh random variables,
even if it revisits an original node. Thus each next-action distribution
conditional on the observed path matches the original process. The
case $\beta=1$ terminates immediately under both representations.
\end{proof}

The common horizon couples hypothetical walks from different seeds
within one sample. Such hypothetical outcomes need not be independent:
only one seed is assigned to any given coupon. The coupling does not
change any single-seed walk law, and independent copies preserve the
independence of the actual coupons. Since $H$ is finite almost surely
and $\E[H+1]=1/\beta$, this is an exact representation, not a
fixed-length truncation. A directed cycle in $G$ becomes a path through
different layers; repeated visits therefore retain fresh actions.

\subsection{Constructing a Reverse Sample}
\label{sec:reverse-construction}

Fix a target user $u$ and one layered realization. Define $R(u)$ as
the set of seeds $s$ whose paths from $(s,0)$ terminate by consumption
at some copy $(u,\ell)$, $0\leq\ell\leq H$. To construct $R(u)$,
let $C_\ell$ be the original nodes whose copies at layer $\ell$ lead
to such a consumption. At the terminal layer,
\[
  C_H=\begin{cases}\{u\},&\text{if $(u,H)$ consumes},\\
                     \varnothing,&\text{otherwise.}\end{cases}
\]
For $\ell=H-1,\ldots,0$, the backward recurrence is
\begin{equation}
\label{eq:backward-layers}
\begin{split}
  C_\ell={}&\{u:\text{$(u,\ell)$ consumes}\}\\
  &{}\cup\{w:\text{$(w,\ell)$ transfers to some }v\in C_{\ell+1}\}.
\end{split}
\end{equation}
The desired reverse set is $R(u)=C_0$.
For each layer, only the root $u$ and in-neighbors of nodes in
$C_{\ell+1}$ need their actions sampled. If an in-neighbor is examined
more than once in that layer, its sampled action is reused; at a
different layer, it receives a fresh draw. Thus mutually exclusive
outgoing choices at one state are respected without fixing an action
for all visits to an original node.

\begin{algorithm}[t]
\caption{Layered Reverse Sample $\textsc{ReverseCoupon}(u)$}
\label{alg:rrset}
\begin{algorithmic}[1]
\Require Root $u$; incoming adjacency lists; $\beta$ and action probabilities
\Ensure Reverse set $R(u)$
\If{$\beta=1$}
  \State $H\gets0$
\Else
  \State Draw $H$ using Eq.~\eqref{eq:random-horizon}
\EndIf
\State $C\gets\{u\}$ with probability $p_u^a/c_u$; otherwise $C\gets\varnothing$
\For{$\ell=H-1,\ldots,0$}
  \State Start an empty action cache for layer $\ell$; $D\gets\varnothing$
  \State Sample and cache the residual action $a_\ell(u)$
  \If{$a_\ell(u)=\mathtt{consume}$}
    \State $D\gets\{u\}$
  \EndIf
  \For{each $v\in C$}
    \For{each $w\in N^-(v)$}
      \If{$a_\ell(w)$ has not been sampled}
        \State Sample and cache it using Eq.~\eqref{eq:residual-actions}
      \EndIf
      \If{$a_\ell(w)=\mathtt{transfer}(v)$}
        \State $D\gets D\cup\{w\}$
      \EndIf
    \EndFor
  \EndFor
  \State $C\gets D$
\EndFor
\State \Return $C$
\end{algorithmic}
\end{algorithm}

An empty $C$ at an intermediate layer is not a stopping condition:
the root may consume at an earlier layer. Only layer zero completes
the sample. In particular, a single failed consumption gate at the
root does not eliminate later or earlier consumption opportunities.

\begin{lemma}[Single-coupon coverage]
\label{lem:coverage}
For every root $u$ and seed $s$,
\begin{equation}
\label{eq:single-reverse-law}
  \Pr[s\in R(u)]=M_{u,s}.
\end{equation}
\end{lemma}

\begin{proof}
At layer $H$, membership in $C_H$ is exactly consumption at $u$.
Inductively, a state at layer $\ell$ leads to consumption at $u$
if it consumes there immediately, or transfers to a state with this
property at layer $\ell+1$. This is precisely
Eq.~\eqref{eq:backward-layers}. Hence $s\in C_0$ if and only if
the layered walk from $s$ is consumed by $u$. Apply
Lemma~\ref{lem:layered-law} and Eq.~\eqref{eq:consume-prob}.
\end{proof}

\subsection{Independent Coupons and Marginal Coverage}
\label{sec:marginal-coverage}

For a fixed labeling $S=\{s_1,\ldots,s_i\}$, generate independent
reverse sets $R_1(u),\ldots,R_i(u)$, each with its own horizon and
state actions. Lemma~\ref{lem:coverage} gives
\begin{equation}
\label{eq:joint-reverse-law}
  \Pr\!\left[\bigvee_{j=1}^{i}s_j\in R_j(u)\right]
       =1-\prod_{s\in S}(1-M_{u,s}).
\end{equation}
Thus, for a uniformly sampled root $U$,
\[
  \sigma(S)=n\,\E\!\left[
    \I\!\left[\bigvee_{j=1}^{i}s_j\in R_j(U)\right]\right].
\]
The correspondence between each seed and its own independent reverse
set is essential. A single shared $R(u)$ has correlated memberships
across candidate seeds, so testing $S\cap R(u)\ne\varnothing$ does
not in general give Eq.~\eqref{eq:joint-reverse-law}.

For greedy selection, we directly estimate marginal gains. Given the
current seed set $S$, sample $U$ uniformly. Generate an independent
reverse set for each seed $s_j\in S$ and test $s_j\in R_j(U)$.
If any test succeeds, return an empty sample. Otherwise generate a
fresh independent reverse set $R_0(U)$ and return
\begin{equation}
\label{eq:marginal-reverse-set}
  R^S=R_0(U)\setminus S.
\end{equation}
The empty sample is retained in the sample count. These samples are
not conditioned on the event that earlier coupons fail to consume
at the root.

\begin{lemma}[Unbiased marginal estimator]
\label{lem:unbiased}
For every fixed $S\subseteq V$ and $v\notin S$,
\begin{equation}
\label{eq:marginal-unbiased}
  n\Pr[v\in R^S]
   =\sum_{u\in V}M_{u,v}\prod_{s\in S}(1-M_{u,s})
   =\Delta(v\mid S).
\end{equation}
\end{lemma}

\begin{proof}
Condition on $U=u$. The selected seeds fail their respective tests
with probability $\prod_{s\in S}(1-M_{u,s})$.
Independently, $v\in R_0(u)$ has probability $M_{u,v}$.
Multiply these probabilities and average over the uniform root.
The last equality is Eq.~\eqref{eq:marginal}.
\end{proof}

This construction separates the two uses of coupling. Within
$R_0(U)$, correlated candidate memberships allow one reverse traversal
to evaluate many alternative seeds. Across actual coupons, independent
reverse sets preserve their independent consumption events.

\subsection{Greedy Selection with Fresh Reverse Samples}
\label{sec:reverse-greedy}

Let $S_a$ be the $k$ users with the largest $p_u^a$, and define
\[
  \mathrm{LB}=\sum_{u\in S_a}p_u^a
       \leq\sigma(S_a)\leq\mathrm{OPT},\qquad
  \mathrm{OPT}=\sigma(S^*).
\]
The first inequality counts immediate consumption at the distinct
seed users. If $\mathrm{LB}=0$ and $k\geq1$, every consumption
probability is zero, and any $k$ distinct users are optimal.
Otherwise, for accuracy $0<\varepsilon<1-1/e$ and failure probability
$0<\delta<1$, set the number of samples per greedy round to
\begin{equation}
\label{eq:sample-budget}
  \theta=\left\lceil
     \frac{10nk^2}{\varepsilon^2\mathrm{LB}^2}
     \ln\frac{2nk}{\delta}
  \right\rceil.
\end{equation}
For the current $S$, generate $\theta$ independent marginal reverse
sets $R_1^S,\ldots,R_\theta^S$ and estimate
\[
  \widehat\Delta(v\mid S)
     =\frac n\theta\sum_{t=1}^{\theta}\I[v\in R_t^S].
\]
Select an unselected user maximizing this estimate. After adding it,
draw new samples for the next round; samples constructed for the old
$S$ are not treated as samples for its enlarged successor.

\begin{algorithm}[t]
\caption{\CIMRIS{}$(G,k,\varepsilon,\delta)$ by Marginal Reverse Sampling}
\label{alg:cim-ris}
\begin{algorithmic}[1]
\Require $1\leq k\leq n$; $\beta>0$; $0<\varepsilon<1-1/e$; $0<\delta<1$
\State Compute $S_a$ and $\mathrm{LB}$
\If{$\mathrm{LB}=0$}
  \State \Return $S_a$
\EndIf
\State Set $\theta$ using Eq.~\eqref{eq:sample-budget}; $S\gets\varnothing$
\For{$i=1,\ldots,k$}
  \State $g_v\gets0$ for every $v\in V\setminus S$
  \For{$t=1,\ldots,\theta$}
    \State Draw $U$ uniformly from $V$; $\mathrm{hit}\gets\mathrm{false}$
    \For{each $s\in S$}
      \State Independently generate $R\gets\textsc{ReverseCoupon}(U)$
      \If{$s\in R$}
        \State $\mathrm{hit}\gets\mathrm{true}$; \textbf{break}
      \EndIf
    \EndFor
    \If{$\mathrm{hit}=\mathrm{false}$}
      \State Independently generate $R\gets\textsc{ReverseCoupon}(U)$
      \For{each $v\in R\setminus S$}
        \State $g_v\gets g_v+1$
      \EndFor
    \EndIf
  \EndFor
  \State Select $v^*\in\arg\max_{v\in V\setminus S}g_v$
  \State $S\gets S\cup\{v^*\}$
\EndFor
\State \Return $S$
\end{algorithmic}
\end{algorithm}

Each round is a coverage-frequency computation over the candidates.
The selected set contains no repeated user. Because the sampled
marginal objective depends on $S$, the analysis below controls errors
round by round rather than applying a fixed-sample maximum-coverage
bound to all rounds.

\subsection{Accuracy and Computational Cost}

\begin{theorem}
\label{thm:approx}
Under Section~\ref{sec:model} with $\beta>0$,
Algorithm~\ref{alg:cim-ris} returns $k$ distinct seeds $\hat S$ and,
with probability at least $1-\delta$, satisfies
\[
  \sigma(\hat S)\geq(1-1/e-\varepsilon)\mathrm{OPT}.
\]
\end{theorem}

\begin{proof}
The zero-spread case is immediate. Otherwise put
$\xi=\varepsilon\mathrm{LB}/(2k)<1/2$.
Condition on the complete history before a round, so that its current
$S$ is fixed. For a candidate $v\notin S$, the fresh observations
$X_t(v)=n\I[v\in R_t^S]$ are independent across $t$, with mean
$\Delta(v\mid S)$. Since an additional coupon activates at most one
new user, $0\leq\Delta(v\mid S)\leq1$. Consequently,
$\operatorname{Var}(X_t(v))\leq n\Delta(v\mid S)\leq n$ and
$|X_t(v)-\E X_t(v)|\leq n$. Bernstein's inequality gives
\[
  \Pr\!\left[|\widehat\Delta(v\mid S)-\Delta(v\mid S)|\geq\xi
      \mid\text{history}\right]
  \leq2\exp\!\left(-\frac{\theta\xi^2}{2n+(2/3)n\xi}\right)
  \leq\frac{\delta}{nk}.
\]
The last step uses Eq.~\eqref{eq:sample-budget} and $\xi<1/2$.
A union bound over at most $n$ candidates and $k$ rounds, conditioning
on each preceding history, bounds the total failure probability by
$\delta$. Dependence between different candidates within a sample
does not affect this argument.

On the resulting event, the selected candidate has true marginal gain
at least $\max_{v\notin S}\Delta(v\mid S)-2\xi$.
By monotonicity and submodularity,
$\max_{v\notin S}\Delta(v\mid S)\geq
(\mathrm{OPT}-\sigma(S))/k$. Writing $F_i=\sigma(S_i)$ for the
spread after $i$ selections therefore gives
\[
  F_i\geq F_{i-1}+\frac{\mathrm{OPT}-F_{i-1}}k-2\xi.
\]
Unrolling this recurrence from $F_0=0$ yields
\begin{align*}
  F_k&\geq\left[1-\left(1-\frac1k\right)^k\right]\mathrm{OPT}-2k\xi\\
     &\geq(1-1/e)\mathrm{OPT}-\varepsilon\mathrm{LB}\\
     &\geq(1-1/e-\varepsilon)\mathrm{OPT}.
\end{align*}
\end{proof}

\paragraph{Sampling cost.}
Precompute incoming adjacency lists and cumulative residual-action
probabilities in $O(n+m)$ time. Conditional on $H$, one backward
layer examines at most $m$ incoming edges and samples at most $n$
state actions. With binary search for categorical actions, this costs
$O(m+n\log(n+1))$ per layer. Since $\E[H+1]=1/\beta$, one reverse
sample has expected cost
\[
  O\!\left(\frac{m+n\log(n+1)}\beta\right).
\]
Only two consecutive active layers and one layer's action cache are
needed, so the sampler uses $O(n+m)$ working storage.

In round $i$, one marginal sample invokes the reverse sampler at most
$i$ times: up to $i-1$ tests for previously selected seeds and one
sample for the candidate coupon. Summing over all rounds gives at
most $\theta k(k+1)/2$ calls. Samples can be processed immediately
into the counters $g_v$ and then discarded. Including sorting for
$S_a$ and scanning candidates, a conservative total expected bound is
\[
  O\!\left(n\log(n+1)+m+nk+
       \frac{\theta k^2\bigl(m+n\log(n+1)\bigr)}\beta\right),
\]
with $O(n+m)$ storage. This bound establishes a finite reverse-sampling
algorithm and its approximation guarantee. It does not imply
near-linear running time: the dependence on $\theta$, $k$, and
$1/\beta$ can be large. Tighter sample budgets and more efficient
reuse of reverse information require additional analysis.

